\documentclass[11pt]{article}
\usepackage[a4paper, top=18mm, bottom=18mm, left=20mm, right=20mm]{geometry}
\usepackage[T2A]{fontenc}
\usepackage[utf8]{inputenc}
\usepackage[ukrainian]{babel}
\usepackage{graphicx}
\setlength{\parindent}{0pt}
\setlength{\parskip}{4pt}
\pagestyle{empty}
\begin{document}
%begin
{\large\textbf{mixed-20}}

\textbf{Варіант \No\,1}\hfill \textit{ПІБ:}\,\underline{\hspace{60mm}}
\vspace{4mm}

%beginitem
1. Що буде виведено за виконання фрагменту коду?
Записати ВИРАЗ, значенням якого є список, що складається з квадратівцілих чисел від 13 до 2003 (включно), що не діляться на 5, записаних в оберненому порядку.
%enditem
%beginitem
2. Відмітити все, що є коректним оператором С++:
( 1 ) cout>>3;

( 2 ) int x+y=z;

( 3 ) cout<<3.4;

( 4 ) int f(int n);

%enditem
%beginitem
3. Що буде виведено за виконання фрагменту коду?

%code
print('R', end=' ')
for c in range(-7,-7,2):
    if c>5:
        continue
        print(c, end=' ')
    if c>5:
        break
    else:
        print(c, end=' ')
    print(c, end=' ')
%code

%enditem
%beginitem
4. Відмітити всі правильні варіанти:
( 1 ) 8,3

( 2 ) x*y

( 3 ) 0X2b

( 4 ) cpp

%enditem
%beginitem
5. Відмітити все, що є коректним оператором С++:
( 1 ) x+3=5

( 2 ) n=2; n--

( 3 ) i,j=j,j

( 4 ) x=1//3

%enditem
%beginitem
6. Що буде виведено за виконання фрагменту коду?
%code

x,y,z=2,3,0
z,x,x,z=x,8,z,y
print(x,y,z)
%code

%enditem
%beginitem
7. Що буде виведено за виконання фрагменту коду?
%code

a,b,c=8,9,6
def f(a,b=7,c):
    print(a,b,c,end="")

f(2,a=3)
print(a,b,c)

%code

%enditem
%beginitem
8. Відмітити все, що є коректним оператором С++:
( 1 ) x+3=5

( 2 ) n=2; n--

( 3 ) i,j=j,j

( 4 ) x=1//3

%enditem
%beginitem
9. Відмітити все, що є коректним оператором С++:
( 1 ) x+3=5

( 2 ) n=2; n--

( 3 ) i,j=j,j

( 4 ) x=1//3

%enditem
%beginitem
10. Що буде виведено за виконання фрагменту коду?
Записати ВИРАЗ, значенням якого є список, що складається з квадратівцілих чисел від 13 до 2003 (включно), що не діляться на 5, записаних в оберненому порядку.
%enditem
%beginitem
11. Що буде виведено за виконання фрагменту коду?
��i���� ��i���� ���� bool* ��� ���������: bool *c, &x, **w, **q, *e, &u;
%enditem
%beginitem
12. Відмітити все, що є коректним оператором С++:
( 1 ) x+3=5

( 2 ) n=2; n--

( 3 ) i,j=j,j

( 4 ) x=1//3

%enditem
%beginitem
13. Відмітити все, що є коректним оператором С++:
( 1 ) cout>>3;

( 2 ) int x+y=z;

( 3 ) cout<<3.4;

( 4 ) int f(int n);

%enditem
%beginitem
14. Що буде виведено за виконання фрагменту коду?
Записати ВИРАЗ, значенням якого є множина, що складається з квадратівцілих чисел від 13 до 2003 (включно), що не діляться на 5.

%enditem
%beginitem
15. Що буде виведено за виконання фрагменту коду?
%code

try:
    print(2,end="")
    print(int("9"),end="")
    print(0,end="")
except ValueError: print(5,end="")
except TypeError: print(4,end="")
else: print(6,end="")
finally: print(7,end="")

%code

%enditem
%beginitem
16. Відмітити всі правильні варіанти:
( 1 ) 8,3

( 2 ) x*y

( 3 ) 0X2b

( 4 ) cpp

%enditem
%beginitem
17. Що буде виведено за виконання фрагменту коду?
%code

a,b,c=8,9,6
def f(a,b=7,c):
    print(a,b,c,end="")

f(2,a=3)
print(a,b,c)

%code

%enditem
%beginitem
18. Що буде виведено за виконання фрагменту коду?
%code
def f():
    try:
        res = int("2")
return 42    except ValueError: return 9
    except TypeError: return 0
else: return 32finally: return 22    return res

print(f())
%code

%enditem
%beginitem
19. Відмітити все, що є коректним оператором С++:
( 1 ) x+3=5

( 2 ) n=2; n--

( 3 ) i,j=j,j

( 4 ) x=1//3

%enditem
%beginitem
20. Відмітити все, що є коректним оператором С++:
( 1 ) x+3=5

( 2 ) n=2; n--

( 3 ) i,j=j,j

( 4 ) x=1//3

%enditem




















%end
\newpage
\end{document}

